\section{Stability, finite charge spaces, and elliptic support}
\label{app:stability-support}

A Hall operation uses families of exact sequences, together with a
coefficient theory on their moduli.  The lattice calculations determine
neither of these.  A product stability condition gives the filtration problem on
$S\times E$.  Moduli and descent conditions then determine which
geometric operations can be defined.

\subsection{A product stability condition}

\begin{definition}[Product heart and central charge]
\label{def:liu-product-stability-datum}
\label{def:retained-ap-liu-heart}
Fix a numerical Bridgeland stability condition
$\sigma_S=(\mathcal A_S,Z_S)$ on $D^b\Coh(S)$ with
$Z_S(\Lambda_S^{\mathrm{alg}})\subset\Q+i\Q$.
Choose a degree-one line bundle $L_E$ on $E$.  The
Abramovich--Polishchuk heart \cite{AbramovichPolishchukTStructures} is
\[
\mathcal A_E^{\mathrm{AP}}=
\{M\in D^b\Coh(X):Rp_{S*}(M\otimes p_E^*L_E^n)\in\mathcal A_S
\text{ for }n\gg0\}.
\]
Write
\[
Z_S(Rp_{S*}(M\otimes p_E^*L_E^n))
=a(M)n+b(M)+i(c(M)n+d(M)).
\]
For rational $t>0$, the weak charge and slope are
\[
Z_t=at-d+ict,
\qquad
\nu_t(M)=\begin{cases}(-at+d)/(ct)&c(M)\ne0,\\
 +\infty&c(M)=0.\end{cases}
\]

For nonzero objects, the minimum and maximum slopes of the
Harder--Narasimhan factors define the torsion pair
\[
\mathcal T_t=\{M:\nu_{t,\min}(M)>0\},
\qquad \mathcal F_t=\{M:\nu_{t,\max}(M)\le0\}.
\]
Both subcategories include zero.  Set
\[
\mathcal A_X^{\mathrm{AP/Liu}}
=\mathcal A_E^t=\langle\mathcal T_t,\mathcal F_t[1]\rangle,
\qquad Z_E^{s,t}=cs+b+i(-at+d).
\]
\end{definition}

\begin{proposition}[The product stability condition]
\label{prop:ap-liu-heart-compatibility}
For rational $s,t>0$, the pair
\[
\sigma_{s,t}=(\mathcal A_E^t,Z_E^{s,t})
\]
is a stability condition whose central charge factors through
$\Gamma_X^{\mathrm{alg}}$.  Its support property holds on
\[
\Lambda_S^{\mathrm{alg}}\oplus
\bigl(\Lambda_S^{\mathrm{alg}}/\ker Z_S\bigr),
\qquad (Q,P)\longmapsto(P,\overline Q).
\]
These pairs extend continuously to positive real parameters.
No support property on the larger full paired lattice is asserted.
\end{proposition}

\begin{proof}
The rational central charge has discrete image because it factors
through the finite-rank lattice $\Lambda_S^{\mathrm{alg}}$.
The assertion is Liu's product theorem, including its torsion-pair
construction and support property
\cite[Theorems~4.7 and 5.9]{LiuProductStability}.
\end{proof}

In the coordinates of Chapter~\ref{ch:k3xe-setup}, these coefficients
have a direct interpretation.  Grothendieck--Riemann--Roch gives
\[
v(Rp_{S*}(M\otimes p_E^*L_E^n))=Q(M)+nP(M).
\]
Indeed, $\operatorname{ch}(L_E^n)=1+n\omega_E$ and
$\operatorname{td}(E)=1$, so integration over $E$ selects $Q+nP$.
Consequently
\[
a(M)+ic(M)=Z_S(P(M)),\qquad
b(M)+id(M)=Z_S(Q(M)).
\]
This identifies the polynomial coefficients with the actual numerical
charges.  The product stability theorem does not assert finite type
or properness for every moduli locus considered below.

\subsection{Finite windows and Harder--Narasimhan filtrations}

At finite height, termination of subobject chains can be imposed by
a nonnegative integer length that strictly increases under strict
inclusions.  This gives an elementary noetherianity criterion for the
chosen subcategory, without a global noetherianity assertion for
\(\mathcal A_X^{\mathrm{AP/Liu}}\).

A finite height bound is generally incompatible with closure under
direct sums.  Indeed, for a nonzero object \(A\), the split inclusions
\[
A\subset A^{\oplus2}\subset A^{\oplus3}\subset\cdots
\]
are strict.  An isomorphism-invariant length strictly increasing on
these inclusions takes infinitely many values.  Such an additive
subcategory cannot have finitely many isomorphism classes.  A finite
window therefore retains exact sequences from the ambient heart,
without being an exact category in its own right.

\begin{definition}[Retained noetherian window]
\label{def:retained-noetherian-window}
Fix the heart \(\mathcal A_X^{\mathrm{AP/Liu}}\).  A retained finite
window at height \(R\) is a full subcategory
\[
\mathcal A_{X,R}^{\mathrm{ret}}\subset\mathcal A_X^{\mathrm{AP/Liu}}
\]
containing zero and closed under isomorphisms, with finitely many
isomorphism classes.  Require every ambient subobject and quotient
of an object in the window to lie in the window.  A retained exact
sequence is a short exact sequence in the ambient heart whose three
terms lie in the window.  No closure under direct sums or arbitrary
extensions is assumed.

The window carries an isomorphism-invariant length function
\[
\ell_R:\operatorname{Obj}(\mathcal A_{X,R}^{\mathrm{ret}})
\longrightarrow\Z_{\ge0}
\]
with \(\ell_R(A)=0\) exactly for zero.
Require \(\ell_R(A)<\ell_R(B)\) for each non-isomorphism ambient
monomorphism \(A\hookrightarrow B\) between retained objects.
\end{definition}

\begin{proposition}[Retained finite windows are noetherian]
\label{prop:retained-window-noetherian}
\textnormal{[proved]}
Every retained finite window
\(\mathcal A_{X,R}^{\mathrm{ret}}\) in the sense of
Definition~\ref{def:retained-noetherian-window} satisfies the
ascending chain condition on ambient subobjects of each retained
object.
\end{proposition}

\begin{proof}
Let
\[
A_0\subset A_1\subset A_2\subset\cdots
\]
be a chain of retained subobjects of an object \(A\) in
\(\mathcal A_{X,R}^{\mathrm{ret}}\).  After deleting repetitions, each
inclusion is strict.  The length function then gives a strictly
increasing sequence
\[
\ell_R(A_0)<\ell_R(A_1)<\ell_R(A_2)<\cdots .
\]
All \(A_i\) lie in the finite retained window and all are subobjects of
\(A\), so \(\ell_R(A_i)\le \ell_R(A)\).  A strictly increasing sequence
of non-negative integers bounded above by \(\ell_R(A)\) is finite.
Thus the original ascending chain stabilizes.
\end{proof}

A chain of strict subobjects of \(A\) has at most \(\ell_R(A)\)
inclusions.  Each inclusion increases an integer length by at least
one.  Stability supplies a second ordering, by phase, which selects
the successive quotients of a Harder--Narasimhan filtration.

\begin{definition}[Retained finite Harder--Narasimhan filtration datum]
\label{def:retained-finite-hn-filtration-datum}
Let \(\mathcal A_{X,R}^{\mathrm{ret}}\) be a retained finite window.
For a nonzero object \(M\) of the heart, its stability phase is
\[
\phi_\sigma(M)=\pi^{-1}\arg Z_E^{s,t}(M)\in(0,1].
\]
An object \(M\) is semistable if every nonzero ambient subobject
\(B\subset M\) has \(\phi_\sigma(B)\le\phi_\sigma(M)\).

Let \(\mathcal S_R\) be the nonzero semistable objects in the
window.  Their phases form a finite ordered set.  Choose an
order-preserving identification of that set with a finite subset
\(\Phi_R\subset\Q\), and denote the resulting phase label by
\(\phi_R:\mathcal S_R\to\Phi_R\).  These rational labels encode the order.  They need not equal the
stability phases.

A retained finite Harder--Narasimhan filtration datum assigns to
every \(A\in\mathcal A_{X,R}^{\mathrm{ret}}\) a filtration
\[
0=A_0\subset A_1\subset\cdots\subset A_m=A
\]
in the ambient heart.  Every \(A_i\) and every quotient
\(E_i=A_i/A_{i-1}\) must belong to the window, with
\(E_i\in\mathcal S_R\) and
\[
\phi_R(E_1)>\phi_R(E_2)>\cdots>\phi_R(E_m).
\]
The zero object has the empty filtration.
\end{definition}

\begin{proposition}[Construction of finite HN filtrations]
\label{prop:retained-window-hn-filtration}
\textnormal{[proved]}
Every retained finite window in the heart of \(\sigma_{s,t}\)
admits a datum of
Definition~\ref{def:retained-finite-hn-filtration-datum}.
For each \(A\), the filtration has at most \(\ell_R(A)\) nonzero
quotients.
\end{proposition}

\begin{proof}
For any nonzero quotient \(C\) of \(A\), choose a nonzero subobject
\(B\subset C\) of maximal phase.  Among subobjects of that phase,
choose one with maximal \(\ell_R(B)\).  The choices exist because
all subobjects of \(C\) belong to the finite window, and phase and
length are constant on isomorphism classes.  Every subobject of
\(B\) is also a subobject of \(C\), so \(B\) is semistable.

We claim that every nonzero subobject \(D\subset C/B\) has phase
strictly smaller than \(\phi_\sigma(B)\).  Let \(B'\subset C\)
be its inverse image.  The ambient exact sequence
\[
0\longrightarrow B\longrightarrow B'\longrightarrow D
\longrightarrow0
\]
has all terms in the window.  Additivity of the central charge gives
\(Z_E^{s,t}(B')=Z_E^{s,t}(B)+Z_E^{s,t}(D)\).

For nonzero numbers with distinct arguments in \((0,\pi]\), the
sum has argument strictly between them.  Equal arguments remain equal.
Thus \(\phi_\sigma(D)>\phi_\sigma(B)\) contradicts maximality of
the phase of \(B\).  Equality would give
\(\phi_\sigma(B')=\phi_\sigma(B)\), whereas the strict inclusion
\(B\subset B'\) gives \(\ell_R(B')>\ell_R(B)\).

This contradicts the second maximality condition and proves the
claim.

Start with \(C=A\), retain the chosen \(B\), and repeat the
construction on \(A/B\).  Taking inverse images in \(A\) produces
a chain of ambient subobjects.  Subobject and quotient closure
keeps all its terms and semistable factors inside the window.
The claim makes their phases strictly decrease.  Each nonzero
factor strictly increases the length of the inverse-image
subobject, so after at most \(\ell_R(A)\) steps the quotient is
zero.  The resulting chain is the required filtration.
\end{proof}

\begin{example}[Bounded semisimple windows]
Let \(S_1,\ldots,S_k\) be pairwise nonisomorphic simple objects of
the ambient heart, and let \(N\) be a nonnegative integer.  The
objects
\[
\bigoplus_{i=1}^k S_i^{\oplus n_i},
\qquad n_i\ge0,\qquad \sum_i n_i\le N,
\]
and their isomorphic copies form a retained finite window.
Subobjects and quotients of a semisimple object are semisimple,
and their simple multiplicities are bounded by those of the object.

The composition length is \(\sum_i n_i\), and strictly increases
under strict monomorphisms.  Each \(S_i\) is stable because it has
no proper nonzero subobject.  Grouping the summands of equal phase
and ordering the groups by decreasing phase gives the HN
filtration explicitly.  For \(k=1\), this is the one-factor
filtration of every nonzero object.

For \(N>0\), the window need not contain the direct sum of two
objects of length \(N\).  It therefore illustrates the distinction
between ambient exact sequences and an additive exact category.
In the AP/Liu heart this construction requires a chosen finite
collection of simple objects.  It does not supply compact Hall
moduli for more general branes.
\end{example}

The construction uses only the stability function, finite object
classes, subobject and quotient closure, and the strict length
bound.  If the length is additive on retained short exact
sequences, it also gives
\(\ell_R(A)=\sum_i\ell_R(E_i)\).  Finite-type moduli and extension
correspondences remain geometric constructions beyond this
filtration theorem.

\begin{definition}[Retained bounded HN type datum]
\label{def:retained-bounded-hn-type-datum}
Let \(\mathcal A_{X,R}^{\mathrm{ret}}\) be a retained finite window
with a retained finite Harder--Narasimhan filtration datum.  Its \emph{retained HN type} is
\[
\operatorname{HNType}_R(A)
=\bigl(([E_1],\phi_R(E_1)),\ldots,([E_m],\phi_R(E_m))\bigr)
\]
where \(E_i=A_i/A_{i-1}\) are the retained semistable quotients.
For zero, use the empty word.

A retained bounded HN type datum consists of a finite set
\(\mathfrak T_R\) containing these words and a decoration of
semistable factor classes:
\[
[E]\longmapsto
\bigl(\widehat c_E,[a_E,b_E],P_E,N_E,\mathcal F_E\bigr).
\]
Here \(\widehat c_E\) is the retained numerical class and
\([a_E,b_E]\) its cohomological-amplitude interval.
The symbols \(P_E,N_E,\mathcal F_E\) specify a Hilbert polynomial,
a regularity bound, and a chosen family closed under the required
operations.  This datum is a
finite bound on active-window HN types.  It is not a finite-type
semistable substack, not a quasi-smooth derived enhancement, and not a
compact Hall stage.
\end{definition}

\begin{proposition}[HN type is bounded in a retained active window]
\label{prop:retained-hn-type-bounded}
\textnormal{[proved]}
Let \(\mathcal A_{X,R}^{\mathrm{ret}}\) be a retained finite window
with a retained finite Harder--Narasimhan filtration datum.  Then the
set of retained HN types occurring in
\(\mathcal A_{X,R}^{\mathrm{ret}}\) is finite.  More precisely, if
\[
N_R=\max\{\ell_R(A)\mid A\in
\mathcal A_{X,R}^{\mathrm{ret}}\},
\]
then every retained HN type has length at most \(N_R\).
\end{proposition}

\begin{proof}
The retained window has finitely many isomorphism classes, hence the
set of retained semistable HN factors is finite.  If
\[
0=A_0\subset A_1\subset\cdots\subset A_m=A
\]
is the retained HN filtration of \(A\), then each inclusion
\(A_{i-1}\subset A_i\) is strict unless \(m=0\).  The retained length
rank gives
\[
0=\ell_R(A_0)<\ell_R(A_1)<\cdots<\ell_R(A_m)\le N_R,
\]
so \(m\le N_R\).  Thus every retained HN type is a word of length at
most \(N_R\) in a finite alphabet of semistable factor classes.  The
set of such words is finite, and the strict phase condition selects a
finite subset.  A finite Liu--Hilbert decoration of the factor classes
therefore gives a finite decorated HN type set.
\end{proof}

The bound \(m\le N_R\) controls the number of factors in each HN
word.  The finite factor alphabet and its Liu--Hilbert decorations
then give finitely many decorated words.  Representing their
semistable factors by algebraic stacks requires additional geometric data.  Finiteness of the word
set alone does not supply those families.

\subsection{Moduli hypotheses and derived amplitudes}

The finite-window argument concerns a finite set of isomorphism
classes.  A finite-type moduli stack can instead contain infinitely
many such classes.  To pass from that argument to geometry, one must
specify the relevant families and prove their boundedness.

\begin{definition}[Bounded moduli data]
\label{def:retained-finite-type-semistable-substack-datum}
\label{def:retained-finite-class-bound-datum}
For a finite set $C_R$ of numerical types, choose algebraic substacks
$\mathfrak M^{ss}_{R,c}$ of the stack of objects of
$\mathcal A_E^t$.  Require each to parametrize exactly the selected
$\sigma_{s,t}$-semistable families of type $c$.  Bounded moduli data include finite-type Quot or Postnikov parameter
spaces covering these stacks and finite sets of Hilbert polynomials.
They also include amplitude bounds $[a_c,b_c]$ and regularity bounds
$N_c$ for cohomology sheaves, relative to a fixed ample bundle on $X$.
When a closed locus in a chosen ambient stack is needed, closedness
is a further hypothesis on the inclusion.
\end{definition}

\begin{proposition}[Consequences of the bounded moduli hypotheses]
\label{prop:retained-finite-type-semistable-substacks}
\label{prop:retained-finite-class-bounds}
Suppose the covering maps in the bounded moduli data are smooth and
surjective.  Then every $\mathfrak M^{ss}_{R,c}$ is of finite type.
Its objects satisfy the specified class, polynomial, amplitude and
regularity bounds.  If the chosen inclusions into ambient stacks are closed immersions,
their images are specialization closed.
\end{proposition}

\begin{proof}
Finite type is smooth-local on an algebraic stack.  The bounds hold
on the covering parameter spaces and descend to their represented
families.  Specialization-closedness is a condition on the chosen
locus, independent of the finiteness of its set of numerical types.
\end{proof}

Derived moduli introduce a different issue.  At an object $F$ of the
heart, the tangent complex of the derived stack of perfect complexes
is
\[
\mathbb T_{\Perf(X),F}=R\Hom_X(F,F)[1].
\]
Negative self-extensions vanish in a heart.  Serre duality on the
threefold then bounds this tangent complex in degrees $[-1,2]$.
Its cotangent complex has amplitude $[-2,1]$ in general.  In
particular, scalar automorphisms give a nonzero degree-$-1$ tangent
space and a nonzero degree-one cotangent space.

\begin{definition}[Derived enhancement and reduced chart]
\label{def:retained-quasi-smooth-derived-enhancement-datum}
A derived enhancement of $\mathfrak M^{ss}_{R,c}$ is a derived stack
$\mathfrak M^{\mathrm{der}}_{R,c}$ with that classical truncation and
a comparison with the relevant moduli of perfect complexes.
A quasi-smooth reduced chart means a separately constructed derived
Deligne--Mumford stack $\mathfrak N_{R,c}$ with perfect cotangent
complex in degrees $[-1,0]$.  Its use as a Calabi--Yau chart requires
a specified $(-1)$-shifted symplectic form and a compatible comparison
with the reduced deformation complexes of the objects.
\end{definition}

\begin{proposition}[The scalar obstruction to a two-term chart]
\label{prop:retained-quasi-smooth-derived-enhancements}
The unrigidified stack of nonzero objects cannot have cotangent
amplitude $[-1,0]$ at an object with scalar automorphisms.
A chosen reduced model is quasi-smooth when its cotangent complex
is perfect with that amplitude.  It is $(-1)$-shifted symplectic only
if its closed two-form induces an equivalence
$\mathbb T\simeq\mathbb L[-1]$.
\end{proposition}

\begin{proof}
The infinitesimal scalar automorphism gives a copy of $\C$ in
$H^{-1}(\mathbb T)$, hence in the dual of $H^1(\mathbb L)$.
This excludes amplitude $[-1,0]$ on the unrigidified stack.
The assertions about the reduced chart are the quasi-smoothness and
nondegeneracy criteria for its specified structures.
\end{proof}

The Calabi--Yau trace gives the ambient shifted symplectic form
\cite{PTVV2013}.  Its restriction along an arbitrary substack need
not be nondegenerate.  Removing scalar automorphisms changes one end
of the tangent complex.  A two-term self-dual model must also account
for the Serre-dual trace direction.  Neither operation is supplied by
boundedness of the classical moduli stack.

\subsection{Scalar and translation quotients}

\begin{definition}[Scalar rigidification]
\label{def:retained-rigidification-inertia-datum}
Suppose the scalar subgroup $\mathbb G_m$ is central in the inertia
of a moduli stack $\mathfrak M$.  A scalar rigidification is a
$\mathbb G_m$-gerbe
\[
\pi:\mathfrak M\longrightarrow\mathfrak M^{\mathrm{rig}}
\]
that replaces every automorphism group by its quotient by the scalar
subgroup.  Require the residual automorphism groups to be finite on
the locus under consideration.
\end{definition}

\begin{proposition}[Residual inertia]
\label{prop:retained-finite-residual-inertia-after-rigidification}
At an object $F$, scalar rigidification has residual stabilizer
$\Aut(F)/\mathbb G_m$.  It is finite precisely when this quotient
is finite.  In particular, removing scalars from
$\bigoplus_i L_i\otimes W_i$ generally leaves a positive-dimensional
group.
\end{proposition}

\begin{proof}
The first assertion is the inertia property of rigidification.
For the displayed polystable sheaf, distinct stable summands have no
Hom groups between them.  Thus
$\Aut(F)=\prod_i\GL(W_i)$, whose quotient by the diagonal scalars
has dimension $\sum_i(\dim W_i)^2-1$ when $F\ne0$.
\end{proof}

Rigidification removes an existing central subgroup of automorphisms.
It is not an external quotient $[\mathfrak M/\mathbb G_m]$.
For example, $B\mathbb G_m\to\operatorname{Spec}\C$ is the scalar
rigidification of $B\mathbb G_m$.  The quotient for a trivial external
$\mathbb G_m$ action instead has another $B\mathbb G_m$ factor.

\begin{definition}[Normalized universal complex]
\label{def:retained-universal-perfect-complex-datum}
Let $\mathcal U\in\Perf(X\times\mathfrak M)$ be the tautological
complex, on which scalar inertia has weight one.  A normalization is
a line bundle $\mathcal L$ on $\mathfrak M$ of scalar weight $-1$.
The normalized complex is
$\mathcal U\otimes\operatorname{pr}_{\mathfrak M}^*\mathcal L$.
If no such line exists, one may instead retain the corresponding
twisted family on the gerbe, with its twist specified.
\end{definition}

\begin{proposition}[The scalar weight of a universal family]
\label{prop:retained-universal-perfect-complexes}
A nonzero tautological complex of scalar weight one does not descend
unchanged across $\pi$.  If a normalizing line exists, the
weight-zero tensor product descends to a perfect complex
$\mathcal U^{\mathrm{rig}}$ on $X\times\mathfrak M^{\mathrm{rig}}$.
This descent is unique up to equivalence and is compatible with
base change.
\end{proposition}

\begin{proof}
The pullback of any complex on $\mathfrak M^{\mathrm{rig}}$ has
trivial action by the kernel $\mathbb G_m$ of the inertia map.
A nonzero weight-one complex has a different action, so it cannot
be such a pullback.

Tensoring by a weight-$-1$ line produces weight zero.
Locally for the faithfully flat topology, the gerbe is
$B\mathbb G_m\times U\to U$.  A weight-zero complex on this product
is precisely the pullback of a complex on $U$: the weight-zero
representation has trivial $\mathbb G_m$ action.  These local
pullback equivalences glue, and are compatible with base change.
Perfectness descends in the same topology.
\end{proof}

A fixed Tor-amplitude bound $[a_c,b_c]$ is further data on a
class-$c$ locus.  This means that $\mathcal U\otimes^L M$ has
cohomology sheaves only in these degrees for every quasicoherent
sheaf $M$.  Tensoring by the normalizing line preserves this bound.
The descended complex $\mathcal U^{\mathrm{rig}}$ and all its
derived pullbacks then have Tor amplitude in $[a_c,b_c]$.
Indeed, the bound is local for faithfully flat descent and is
preserved by derived pullback.

On $B\mathbb G_m$, the standard one-dimensional representation is
the simplest example of the obstruction.  Its underlying vector
space is a line, but its inertia action prevents descent to a point.
A normalization is additional data and need not exist globally.
Determinants of cohomology give an explicit sufficient condition.

\begin{proposition}[Normalization from Euler weights]
\label{k3:prop:scalar-normalization}
Let $\mathcal U$ be a perfect family of scalar weight one on
$X\times\mathfrak M$, where $X$ is smooth and projective.  Suppose
there are perfect complexes $A_1,\ldots,A_k$ on $X$ for which the
integers
\[
d_j=\chi(X,\mathcal U_m\otimes^L A_j)
\]
are independent of $m\in\mathfrak M$ and have greatest common
divisor one.  Choose integers $b_j$ with $\sum_jb_jd_j=-1$.
For the projection $p:X\times\mathfrak M\to\mathfrak M$, the line
\[
\mathcal L=\bigotimes_j
\left(\det Rp_*(\mathcal U\otimes^L\operatorname{pr}_X^*A_j)
\right)^{b_j}
\]
has scalar weight $-1$ and normalizes the universal family.
\end{proposition}

\begin{proof}
Proper smooth pushforward makes each displayed complex perfect.
A scalar $z$ acts by $z$ on every cohomology group of its fiber.
It therefore acts by $z^{d_j}$ on its determinant, since the
exponent is the alternating sum of the cohomology dimensions.
The tensor product has weight $\sum_jb_jd_j=-1$.
Proposition~\ref{prop:retained-universal-perfect-complexes} then
gives the descended normalized family.
\end{proof}

For a family of virtual rank one, take $A_1=\mathcal O_x$ at any
fixed smooth point $x\in X$.  Derived restriction has Euler
characteristic one, so the inverse of its determinant is a
normalizing line.  This provides a normalization on a rank-one
moduli locus without a choice of a frame of the universal family.

\begin{definition}[Elliptic quotient and finite closed cover]
\label{def:retained-e-translation-rigidification-datum}
\label{def:retained-closed-substack-datum}
Suppose translation in $E$ acts on $\mathfrak M^{\mathrm{rig}}$ and
preserves the selected locus.  Let
$\mathfrak Q=[\mathfrak M^{\mathrm{rig}}/E]$.
When a free quotient is used, require the action to be free in the
stack-theoretic sense.  When finite stabilizers are allowed, specify
them on the quotient.  A finite closed cover consists of closed
immersions $\mathfrak Z_j\hookrightarrow\mathfrak Q$ whose images
cover that locus, with finite inertia on each $\mathfrak Z_j$.
\end{definition}

\begin{proposition}[Quotient and covering consequences]
\label{prop:retained-e-translation-rigidifications}
\label{prop:retained-closed-substacks}
The map $\mathfrak M^{\mathrm{rig}}\to\mathfrak Q$ is an
$E$-torsor in the stack sense.  The $\mathfrak Z_j$ form the stated
finite closed cover with the prescribed inertia.  Equivariant coefficients descend only with
a cocycle compatible with this quotient and its stabilizers.
\end{proposition}

\begin{proof}
A map $T\to[\mathfrak M^{\mathrm{rig}}/E]$ specifies an $E$-torsor
on $T$ and an equivariant map from that torsor to
$\mathfrak M^{\mathrm{rig}}$.  This is the defining pullback property
of the quotient.  The covering and inertia conclusions are precisely
the hypotheses on the closed immersions.  Descent of coefficients
uses their equivariant cocycle, which is separate from the existence
of the quotient stack.
\end{proof}

\subsection{A finite word model for closure}

The following calculation records three closure conditions without
asserting that a geometric correspondence has been constructed.
Let $s_3,s_2,s_1$ be symbols with strictly decreasing phases and set
\[
\mathcal W=\{0,s_3,s_2,s_1,e_{32},e_{21},f_{321}\},
\quad e_{32}=(s_3,s_2),\quad e_{21}=(s_2,s_1),
\quad f_{321}=(s_3,s_2,s_1).
\]

\begin{definition}[Selected extensions and factor closure]
\label{def:retained-extension-closure-datum}
\label{def:retained-hn-factor-closure-datum}
Retain the four concatenation rules
\[
s_3*s_2=e_{32},\quad s_2*s_1=e_{21},\quad
 e_{32}*s_1=f_{321},\quad s_3*e_{21}=f_{321}.
\]
For a geometric interpretation, each rule must be represented by an
actual exact sequence with the indicated types.  Factor closure
requires a chosen semistable moduli locus for every factor occurring
in a retained word.  Each factor retains its numerical class and,
when a finite closed cover is used, its chosen closed locus.
Require every factor occurrence to belong to these loci.
\end{definition}

\begin{proposition}[Closure of the word set]
\label{prop:retained-extension-closure}
\label{prop:retained-hn-factor-closure}
The four selected concatenations lie in $\mathcal W$.  Every nonzero
factor is one of $s_3,s_2,s_1$, and the six nonempty words have ten
factor occurrences in total.  Under the stated factor-closure
requirements, taking these factors stays in the chosen substack
system.
\end{proposition}

\begin{proof}
The two words of length two contribute four occurrences, the three
words of length one contribute three, and $f_{321}$ contributes three.
The displayed rules list their retained concatenations.
\end{proof}

\begin{definition}[Duality on words]
\label{def:retained-dual-closure-datum}
Set $D(s_3)=s_1$, $D(s_2)=s_2$, $D(s_1)=s_3$ and
$D(t_1,\ldots,t_k)=(D(t_k),\ldots,D(t_1))$.  Set $D(0)=0$.
\end{definition}

\begin{proposition}[Dual closure]
\label{prop:retained-dual-closure}
The map $D$ is an involution of $\mathcal W$.  It exchanges
$e_{32}$ and $e_{21}$ and fixes $f_{321}$.
\end{proposition}

\begin{proof}
Reversing the word twice restores its order, and dualizing each
symbol twice restores the symbol.  Applying the rule to the three
longer words gives the asserted values.
\end{proof}

These statements concern selected words and selected extensions.
They assert neither closure under every extension nor existence of
an extension stack.  A sheaf duality must preserve the selected heart, or be compared
with a second heart.  It must also act on families and extension flags.

\subsection{Effective charges and parameter transport}

\begin{definition}[Geometric charge support]
\label{def:algebraic-effective-hall-support}
Choose an extension-closed full subcategory $\mathcal B$ of the heart
$\mathcal A_E^t$, containing zero.  Its effective charge semigroup is
\[
\Gamma_{X,\mathcal B}^{\mathrm{eff}}
=\{v_X(F):F\in\mathcal B\}\subset\Gamma_X^{\mathrm{alg}}.
\]
For a fixed phase $\phi$, the semistable subcategory
$\mathcal P(\phi)$ is one such choice.  The notation
$\Gamma_{X,\sigma}^{\mathrm{eff}}$ presupposes a specified heart or
single-phase subcategory of this kind.  A bounded window instead
has a charge set, without an assertion of additive closure.
\end{definition}

The semigroup property follows from direct sums, which are split
extensions.  Semistable objects of different phases need not have a
semistable direct sum.  Their union is a support set and is not
assumed to define a semigroup.

\begin{remark}[Formal, algebraic and effective charges]
\label{rem:three-tier-inclusion}
The inclusions
\[
\Gamma_{X,\mathcal B}^{\mathrm{eff}}
\subset\Gamma_X^{\mathrm{alg}}
\subset\Gamma_X^{\mathrm{form}}
\]
separate existence of objects from algebraicity of numerical classes
and from cohomological lattice arithmetic.  Only the last two groups
are fixed by their lattice definitions.  The first depends on the
chosen category and stability condition.
\end{remark}

\begin{definition}[Admissible parameter transport]
\label{def:admissible-dirac-igusa-parameter-datum}
A parameter datum is
\[
\mathfrak p=(S,H,\sigma,E,0_E,\tau_X,\mathcal R_{\mathrm{HN}}),
\]
where $\tau_X$ trivializes $K_X$ and $\mathcal R_{\mathrm{HN}}$ is a
specified cofiltered system of finite stages.  Every stage includes
those moduli, correspondences, coefficients and orientation data that
its operations use.  An admissible isomorphism includes maps $f:S\to S'$, $u:E\to E'$
preserving these structures.  It also specifies compatible
isomorphisms of the finite-stage data after a common cofinal refinement.  Write
$\mathsf{Par}^{\mathrm{DI}}$ for this groupoid.
\end{definition}

\begin{proposition}[Transport of a specified system]
\label{prop:admissible-base-change-functoriality}
An admissible isomorphism transports the finite-stage constructions
and induces their pro-isomorphism.  It preserves the Mukai pairing,
$\Pi_X$, $B$, and $\overline\Pi_X$.  It preserves Hall, Koszul,
primitive and Pfaffian structures when their comparison maps and
coherence identities are part of the given stage isomorphisms.
\end{proposition}

\begin{proof}
The map $f$ induces a Mukai isometry, while $u$ preserves the two
integral generators of $H^{\mathrm{even}}(E)$.  The formulas for the
charge maps are therefore invariant.  Derived transport along
$f\times u$ carries the given families and correspondences to their
primed versions.  The supplied coefficient and orientation
isomorphisms identify the operations.  Compatibility with transitions
then gives an isomorphism of the cofiltered systems after their
common cofinal refinement.
\end{proof}

\begin{definition}[Polarization and origin]
\label{def:polarization-origin-dependence}
Changing $(H,\sigma)$ changes the selected moduli and their
Harder--Narasimhan types.  A comparison across stability chambers
requires specified wall-crossing maps with their coefficient and
orientation compatibilities.  For the origin change from $0_E$ to
$a\in E$, define the determinant anchor by
\[
\lambda_a^{\det}(\mathcal F)=
\det Rp_{E*}\mathcal F\otimes\mathcal O_E(-\chi(\mathcal F)a),
\]
for a perfect complex on $X$.  Here
$\chi(\mathcal F)=\chi(X,\mathcal F)$.
\end{definition}

\begin{proposition}[Origin change of the determinant anchor]
\label{prop:elliptic-origin-independence-character}
The anchors satisfy the Picard-group identity
\[
[\lambda_a^{\det}(\mathcal F)]-
[\lambda_0^{\det}(\mathcal F)]
=[\mathcal O_E(\chi(\mathcal F)(0_E-a))].
\]
This identity does not assign a character to a finite stabilizer.
If an equivariant line and compatible transport of all operations are
specified, their chosen character twist is part of that transport.
\end{proposition}

\begin{proof}
Canceling the two copies of $\det Rp_{E*}\mathcal F$ gives the
line-bundle ratio.  A stabilizer character requires an action on a
fiber of a line bundle, hence a linearization.  The underlying line does not determine it.  On a point with
trivial action by a finite group $H$, every character of $H$
linearizes the trivial line.  Thus a Picard class alone cannot specify the
character needed for equivariant transport.
\end{proof}

Even after the anchor ratio is fixed, orientation lines can carry
independent characters and higher compatibility conditions.
A coherent equivariant comparison must supply these explicitly.
The determinant translation law below describes the numerical
part of this problem.

\subsection{Finite regrading and radicals}

\begin{remark}[Finite charge spaces]
\label{rem:finite-charge-window-verification}
Choose a finite set of source charges together with the maps whose
images are needed for a given calculation.  Closure under those maps
is part of the choice.  A finite charge set, an HN type bound and a
finite-type moduli stack are different objects.  A denominator
identity supplies none of their transition maps.
\end{remark}

\begin{definition}[Gram fibers and parity]
\label{def:gram-fibres-zero-fibre}
Write $\Gamma_X(\gamma)=\Pi_X^{-1}(\gamma)$ and
\[
K_{\Pi,0}=\Gamma_X(0)=\{(Q,P):Q^2=P^2=Q\cdot P=0\}.
\]
For a finite-dimensional super vector space $V=\bigoplus_cV_c$, set
\[
(\Pi_{X,*}^{\mathrm{fib}}V)_{\gamma,\bar p}
=\bigoplus_{c\in\Gamma_X(\gamma)}V_{c,\bar p},
\qquad\bar p\in\Z/2.
\]
\end{definition}

\begin{proposition}[The zero Gram fiber]
\label{prop:zero-gram-fibre-not-subgroup}
The set $K_{\Pi,0}$ is not a subgroup of
$\Gamma_X^{\mathrm{form}}$.
\end{proposition}

\begin{proof}
In two orthogonal hyperbolic planes, set
$k_1=(e_1,e_2)$ and $k_2=(f_1,-f_2)$.  Both have Gram degree zero,
whereas
\[
\Pi_X(k_1+k_2)=\Pi_X(e_1+f_1,e_2-f_2)=(1,0,-1).
\]
\end{proof}

\begin{definition}[A radical in the zero-degree space]
\label{def:zero-fibre-radical-finite-source}
For $V_{\Pi,0}=\bigoplus_{c\in K_{\Pi,0}}V_c$ and a bilinear form
$h$ on this finite-dimensional space, define
\[
\operatorname{Rad}_{\Pi,0}(V,h)
=\ker\bigl(V_{\Pi,0}\longrightarrow V_{\Pi,0}^*\bigr),
\qquad v\longmapsto h(v,-).
\]
\end{definition}

\begin{remark}[Radical and cancellation]
\label{rem:fibre-summation-not-radical}
This radical is a vector subspace.  It gives no quotient of the charge
group by $K_{\Pi,0}$.  Nor does it automatically define an ideal or a
coideal in an algebra carried by $V$.  For example, the form with
matrix $\left(\begin{smallmatrix}0&0\\0&1\end{smallmatrix}\right)$
on $\Q z_0\oplus\Q z_1$ has radical $\Q z_0$.

A different loss of information occurs when one takes superdimension.
An even line and an odd line have parity dimensions $(1,1)$ and
superdimension zero.  Such lines can occur in one Gram degree with
distinct source charges.  For instance,
\[
\Pi_X(e_1+f_1,f_1+e_2)=\Pi_X(e_1+f_1,e_1)=(1,1,0).
\]
Assigning opposite parities to these two summands gives this example.

Exchanging $Q$ and $P$ in
$(e_1+f_1,f_1+e_2+f_2)$ gives another pair, both of degree $(1,1,1)$.
The two zero-degree charges in the previous proof give the same
parity phenomenon inside $V_{\Pi,0}$.
\end{remark}

\begin{proposition}[Exactness and parity of finite regrading]
\label{prop:finite-fibre-pushforward-exact}
\label{prop:finite-parity-pushforward}
Regrading finite-support super vector spaces along any set map is
exact and preserves the even and odd dimensions separately, by
summing them in each fiber.  This applies to $\Pi_X$ and to the
homomorphism $\overline\Pi_X$ on its different source group.
\end{proposition}

\begin{proof}
Every target component is a finite direct sum of source components.
Direct sums preserve exact sequences, with no change of parity.
Additivity of the set map is unnecessary for this vector-space
statement.
\end{proof}

\begin{proposition}[Finite bar and bracket operations]
\label{prop:finite-bar-pushforward-commutes}
\label{prop:finite-bracket-pairing-pushforward}
Let $f=\overline\Pi_X$.  Let $A$ be a $\widehat\Gamma_X$-graded augmented differential
graded algebra with finite support.  Regrading commutes with the
reduced tensor coalgebra on $\bar A[1]$, truncated to lengths
$1$ through $N$.  It preserves deconcatenation and the bar
differential.  Regrading a homogeneous Lie superalgebra preserves
its bracket.  A degree-zero pairing remains of degree zero.
\end{proposition}

\begin{proof}
Proposition~\ref{k3:prop:regrading} gives the tensor and differential
identifications, including the cohomological shift $[1]$.
On a word $[a_1|\cdots|a_j]$, additivity gives target degree
$\sum_i f(\deg a_i)$.  Cutting the word or multiplying adjacent
factors preserves this degree.  The same equality proves bracket
homogeneity.  If $h(V_a,V_{a'})$ is nonzero only for $a+a'=0$,
then it is nonzero only for $f(a)+f(a')=0$ after regrading.
\end{proof}

The unshifted tensor coalgebra has the same degree calculation.
The shift is necessary when it is used as the cohomological bar
complex.  Regrading does not construct a chiral Koszul comparison
or establish a pairing identity that was absent in the source.

\begin{proposition}[Filtrations and height decompositions]
\label{prop:finite-pbw-pushforward}
\label{prop:finite-triangular-pushforward}
Let $V$ be a finite-support $\widehat\Gamma_X$-graded vector space
with a filtration $F_dV$ by homogeneous subspaces.  Then
\[
\operatorname{gr}_d(f_*V)_\gamma
=\bigoplus_{f(a)=\gamma}\operatorname{gr}_dV_a.
\]
A homogeneous filtered multiplication with
$F_iV\,F_jV\subset F_{i+j}V$ retains this property after regrading.
For any homomorphism $h:\Gamma_{\mathrm{gram}}\to\Z$, the direct
sum decomposition by negative, zero and positive values of $h\circ f$
is also preserved.  A bracket has height equal to the sum of its
input heights.
\end{proposition}

\begin{proof}
Exactness commutes the quotient $F_d/F_{d-1}$ with the sum in each
fiber.  Multiplication keeps its filtration bound because the
underlying map is unchanged.  The height sign is constant on a fiber
of $f$, and $h(f(a+a'))=h(f(a))+h(f(a'))$.
\end{proof}

This includes a supplied Poincar\'e--Birkhoff--Witt filtration.
When the filtered pieces have finite dimension, put
$b_d(\gamma)=\dim F_d(f_*V)_\gamma$.  Their graded dimension is
$b_d(\gamma)-b_{d-1}(\gamma)$.

Positive-positive brackets stay
positive, negative-negative brackets stay negative, and zero-height
inputs preserve the other input's height.  For opposite signs, the
actual sum decides the output.  The zero-height space need not be
a Cartan subalgebra.  A completed triangular decomposition requires
its own algebraic and continuity hypotheses.

\subsection{Local and wrapped moduli}
\label{subsec:hybrid-carrier-definition}

\begin{definition}[Hybrid parameter space]
\label{def:hybrid-ran-carrier}
At a finite stage $R$, choose local color classes with projection-finite
moduli $\mathfrak M^{L}_{R,c}$ and wrapped color classes with
prequotient moduli $\mathfrak M^{W}_{R,c}$.  A test family contains a finite configuration of maps to $E$ with
local colors.  It also contains finitely many wrapped families
with their specified anchors.

Maps identify repeated local
points and transport all labels and families coherently.  Denote this
prestack by $\operatorname{Ran}^{\mathrm{hyb}}_{R,\mathrm{pre}}(E)$.
When its descent data for a topology $\tau$ are specified, write
$\operatorname{Ran}^{\mathrm{hyb}}_{R,\tau}(E)$ for the resulting
stack.  Empty configurations are included when a unit is required.
\end{definition}

The local and wrapped subprestacks are denoted
$\operatorname{Ran}^{\mathrm{loc}}_{R,\mathrm{pre}}(E)$ and
$\operatorname{Ran}^{\mathrm{wr}}_{R,\mathrm{pre}}(E)$.
For effective curves, a positive projection degree gives a wrapped
object by Lemma~\ref{lem:projection-locality-obstruction}.
For more general complexes the support condition is part of the
chosen moduli problem.  The degree-labeled set
$\bigsqcup_{b\ge0}\operatorname{Ran}(E)\times\{b\}$ does not retain
these wrapped families.

\begin{remark}[Factorization is additional structure]
\label{rem:two-strata-two-registers}
A projection-finite moduli problem can be organized over the Ran
prestack.  A factorization object on that prestack still requires
restriction, collision and disjoint-support maps with their
compatibilities.  Wrapped families have their own parameter stacks.
The hybrid parameter space records both types and does not itself
construct their operations.
\end{remark}

\subsection{Mixed extension correspondences}
\label{subsec:mixed-correspondences}

\begin{definition}[Selected extension spans]
\label{def:mixed-correspondences}
For $A,B\in\{L,W\}$ and classes $c_1,c_2,c_3$ with
$c_3=c_1+c_2$, specify a stack $\mathfrak E^{AB}_{R,c_1,c_2;c_3}$
of exact sequences $0\to F_1\to F_3\to F_2\to0$ in the chosen
heart.  It has a span
\[
\mathfrak M^A_{R,c_1}\times\mathfrak M^B_{R,c_2}
\xleftarrow{\ p\ }
\mathfrak E^{AB}_{R,c_1,c_2;c_3}
\xrightarrow{\ q\ }\mathfrak M^C_{R,c_3},
\]
where $p$ takes the subobject and quotient, $q$ takes the middle
object, and $C$ is its specified support type.
Each operation requires the pullback or Gysin map along $p$, the
pushforward along $q$, and compatible coefficient maps with their
degree shifts.  In a Borel--Moore theory, properness of $q$ is one
sufficient pushforward hypothesis.
\end{definition}

There are four ordered binary types $LL,LW,WL,WW$.  The two mixed
types have different input orders even when their charge sums agree.
For local configurations indexed by finite sets $I,J$, the target
configuration is the union of the two configurations, with repeated
points identified.  In a mixed span the wrapped source anchor and
the wrapped target anchor are both retained.  In a $WW$ span there
are two wrapped source anchors.  These functions belong to the
prequotient correspondence and require their own compatibility laws.

A local--wrapped--local flag has three source factors and a wrapped
target.  Its two binary decompositions involve the ordered types
$LW$ then $WL$, or $WL$ then $LW$, with the appropriate intermediate
class in each case.  The corresponding flag stacks must identify the
two composites together with their coefficient and orientation maps.
This is the content required of the two-sided mixed correspondence.

If translation acts equivariantly on such a span, quotient all its
terms by the same copy of $E$.  For $E$-stacks $M_1,M_2$, there is
a canonical equivalence
\[
[(M_1\times M_2)/E]
\simeq[M_1/E]\mathbin{\times}_{BE}[M_2/E].
\]
An object on either side is one $E$-torsor with two equivariant maps.
Thus the source of the quotient extension span is this fiber product
over $BE$.  Two unrelated torsors, as in the ordinary product of the
two quotients, do not give the same source.  The empty-object quotient
is $BE$, which is the corresponding relative unit.

\begin{remark}[Flags and higher coherence]
\label{rem:eight-words-vs-colored-trees}
The eight words
\[
LLL,\ LLW,\ LWL,\ LWW,\ WLL,\ WLW,\ WWL,\ WWW
\]
list the ternary support types.
An exact category has canonical comparisons between the two ways of
forming a three-step extension flag.  Restricted moduli and their
coefficients must inherit these comparisons.  Units, permutations,
longer flags, collisions and descent impose further identities.
They are indexed by colored trees and their contractions, not by
the eight ternary words alone.
\end{remark}

\subsection{Determinant anchors and translation}
\label{subsec:e-quotient-anchor}

\begin{definition}[Determinant anchor]
\label{def:determinant-anchor}
Let $\mathcal F$ be a perfect family on $S\times E\times T$, and
let $\pi_{E,T}:S\times E\times T\to E\times T$ be the projection.
Suppose $\chi(\mathcal F_t)=\chi$ is constant on $T$.  Set
\[
\mathscr D(\mathcal F)=\det R\pi_{E,T*}\mathcal F,
\qquad
\lambda^{\det}(\mathcal F)
=\mathscr D(\mathcal F)\otimes
\operatorname{pr}_E^*\mathcal O_E(-\chi\,0_E).
\]
We use $E\simeq\Pic^0(E)$ given by
$a\mapsto\mathcal O_E(a-0_E)$ to regard the relative Picard class
of $\lambda^{\det}$ as a map $T\to E$.
\end{definition}

Proper smooth pushforward makes $R\pi_{E,T*}\mathcal F$ perfect,
so its determinant is defined and compatible with derived base change
\cite{KnudsenMumfordDet,DeligneDetCohomology}.
For a perfect complex $C$ on a genus-one curve,
$\deg\det C=\chi(C)$.  Hence the normalized line has relative degree
zero.  Its class in the relative Picard functor disregards lines
pulled back from $T$.  It does not choose a trivialization of such lines.

\begin{proposition}[Translation weight]
\label{prop:determinant-anchor-translation-weight}
Use support translation
$a\cdot\mathcal F=(\operatorname{id}_S\times\tau_a)_*\mathcal F$.
Then
\[
\lambda^{\det}(a\cdot\mathcal F)
=\lambda^{\det}(\mathcal F)+\chi a
\quad\text{in }\Pic^0(E)\simeq E.
\]
For a section $a:T\to E$, the ratio in the relative Picard functor
is represented by $\mathcal O_{E\times T}(\chi(\Gamma_a-\Gamma_0))$, where
$\Gamma_a$ is its graph.  The kernel of the translation weight is
$E[|\chi|]$ for $\chi\ne0$, and is all of $E$ for $\chi=0$.
\end{proposition}

\begin{proof}
Base change identifies the determinant after translation with
$\tau_{-a}^*\mathscr D(\mathcal F)$.
For a degree-$\chi$ line $L$ on $E$, the relative Picard class of
$\tau_{-a}^*L\otimes L^{-1}$ is $\chi(a-0_E)$.
One sees this first for a divisor by translating each point, then
for any line by subtracting divisors.  The same identity holds in
families in the relative Picard functor.  Normalizing both sides by
$-\chi\,0_E$ gives the result.  The sign is positive for support
pushforward and negative for translation pullback.
\end{proof}

\begin{proposition}[Unit-weight cover and finite slice]
\label{k3:prop:anchor-cover}
Let $\mathfrak M$ be an $E$-stack with a morphism
$\lambda:\mathfrak M\to E$ satisfying
$\lambda(a\cdot F)=\lambda(F)+\chi a$ for a fixed nonzero integer
$\chi$.  Put
\[
\widetilde{\mathfrak M}
=\mathfrak M\mathbin{\times}_{E,[\chi]}E,
\qquad \mathfrak M_0=\lambda^{-1}(0_E).
\]
Then $\widetilde{\mathfrak M}\to\mathfrak M$ is finite \emph{\'etale}
of degree $\chi^2$.  The anchor $(F,b)\mapsto b$ has translation
weight one under $a\cdot(F,b)=(a\cdot F,b+a)$.  Moreover,
\[
\widetilde{\mathfrak M}\simeq E\times\mathfrak M_0,
\qquad
[\mathfrak M/E]\simeq[\mathfrak M_0/E[|\chi|]].
\]
\end{proposition}

\begin{proof}
Multiplication by $\chi$ on a complex elliptic curve is an \emph{\'etale}
isogeny of degree $\chi^2$.  Its base change gives the first assertion.
The equation $\chi b=\lambda(F)$ is preserved by the displayed
$E$ action.  The mutually inverse maps
\[
(F,b)\longmapsto (b,(-b)\cdot F),
\qquad (b,F_0)\longmapsto(b\cdot F_0,b)
\]
identify the cover with $E\times\mathfrak M_0$.

Put $K=E[|\chi|]$.  The isogeny $[\chi]$ is a $K$-torsor, so
$\mathfrak M$ is the quotient of this cover by its deck action.
Under the displayed identification that action is
\[
k\cdot(b,F_0)=(b-k,k\cdot F_0).
\]

It commutes with the $E$ action on the first factor.  Quotienting
$E\times\mathfrak M_0$ first by $E$ leaves $\mathfrak M_0$ with
its $K$ action.  Quotienting in the other order gives
$[\mathfrak M/E]$.  This proves the equivalence of stacks, including
their automorphism groups.
\end{proof}

This constructs the normalization on the nonzero-$\chi$ locus from
the actual anchor map.  Orientation and coefficient descent still
have to be compatible with the finite residual group.  The cover
does not choose its stabilizer characters.

\begin{proposition}[The zero-Euler branch]
\label{prop:determinant-anchor-chi-zero-degeneracy}
When $\chi=0$, the determinant anchor is constant along every
translation orbit.  On a locus where translation acts freely, it
therefore fails to distinguish all points of that one-dimensional
orbit.  An equivariant anchor of weight one cannot be a function
of this determinant anchor alone.
\end{proposition}

\begin{proof}
The translation formula becomes
$\lambda^{\det}(a\cdot F)=\lambda^{\det}(F)$.
A function of this invariant map is invariant, whereas a
weight-one anchor changes by $a$.  These requirements disagree for
$a\ne0$ on a free orbit.
\end{proof}

\begin{proposition}[Jordan--H\"older data beyond the determinant]
\label{prop:determinant-anchor-extra-data-separates-rank-two}
For a semistable rank-$r$, degree-zero vector bundle $V$ on $E$, its
Jordan--H\"older factors are degree-zero lines $L_1,\ldots,L_r$.
The divisor
\[
\operatorname{sp}^{\mathrm{JH}}(V)=\sum_i[L_i]
\in\operatorname{Sym}^r\Pic^0(E)
\]
is independent of the filtration and has Abel sum $\det V$.
It distinguishes $\mathcal O_E^{\oplus2}$ from $L\oplus L^{-1}$
for every nontrivial $L\in\Pic^0(E)$, although their determinants
are equal.
\end{proposition}

\begin{proof}
The classification of semistable bundles on an elliptic curve gives
the stated stable factors and identifies their unordered collection
with the S-equivalence class
\cite{AtiyahVectorBundlesElliptic,TuSemistableElliptic}.
The determinant is the tensor product of these factors, giving the
Abel sum.  If $[L]+[L^{-1}]=2[\mathcal O_E]$, equality of the two
unordered multisets forces $L=\mathcal O_E$.
\end{proof}

This divisor distinguishes the two S-equivalence classes, but not
points of a translation orbit.  Translation acts trivially on the
isomorphism class of each degree-zero line, so it also fixes the
divisor.  The zero-$\chi$ orbit-separation problem remains after
adding this particular invariant.

\subsection{Topological obstruction groups}

We compute the cohomology of classifying spaces with coefficients
$\mathbb F_2$.  Here $BE$ means the topological classifying space of
the complex Lie group $E$.  To use these groups for an algebraic
orientation problem, one must specify the comparison carrying its
actual equivariant obstruction classes into them.

\begin{lemma}[Connected translations]
\label{lem:connected-e-descent}
There is an isomorphism
\[
H^*(BE;\mathbb F_2)=\mathbb F_2[u_1,u_2],\qquad |u_i|=2.
\]
Thus $H^1(BE;\mathbb F_2)=0$, whereas a degree-two class has the
form $\alpha=a_1u_1+a_2u_2$ and vanishes exactly when $a_1=a_2=0$.
\end{lemma}

\begin{proof}
Topologically $E\simeq(S^1)^2$, so
$BE\simeq(\mathbb{CP}^{\infty})^2$.  The two universal first Chern
classes and the K\"unneth formula give the polynomial ring.
\end{proof}

Vanishing of $H^1$ excludes a continuous sign character of the
connected group.  It does not force a degree-two gerbe class to vanish.

\begin{lemma}[Finite translation groups]
\label{lem:finite-stabilizer-e-descent}
For $E[2]\cong(\Z/2)^2$,
\[
H^*(BE[2];\mathbb F_2)=\mathbb F_2[x_1,x_2],\qquad |x_i|=1.
\]
Write a degree-two class as
$\beta=b_{20}x_1^2+b_{11}x_1x_2+b_{02}x_2^2$.
Its restrictions to the cyclic subgroups generated by $(1,0)$,
$(0,1)$ and $(1,1)$ have coefficients $r_1,r_2,r_3$ satisfying
\[
(b_{20},b_{11},b_{02})=(r_1,r_1+r_2+r_3,r_2).
\]
A degree-one class $\lambda=\lambda_1x_1+\lambda_2x_2$ instead has
restriction coefficients $(\lambda_1,\lambda_2,\lambda_1+\lambda_2)$.

If $2^a$ is the exact power of two dividing $N$, with $a\ge2$, then
\[
H^*(BE[N];\mathbb F_2)
=\Lambda(x_1,x_2)\otimes\mathbb F_2[y_1,y_2],
\qquad |x_i|=1,\quad |y_i|=2.
\]
Its degree-two classes are
$A_1y_1+A_{12}x_1x_2+A_2y_2$.
The mixed class $x_1x_2$ is invisible on every cyclic subgroup.
For $N$ odd, all positive-degree mod-two cohomology vanishes.
For $a=1$, the first polynomial-ring description applies.
\end{lemma}

\begin{proof}
The classifying space of $\Z/2$ is $\mathbb{RP}^{\infty}$.
Restriction to the three subgroups substitutes $(x,0)$, $(0,x)$
and $(x,x)$, giving the displayed formulas.

For $\Z/2^a$ with $a\ge2$, the periodic resolution gives one
mod-two cohomology class in each degree, with a periodicity generator
$y$ in degree two.  Its degree-one generator $x$ lifts to a
$\Z/4$ character.  Its Bockstein, hence $x^2$, is zero.
The cyclic ring is therefore $\Lambda(x)\otimes\mathbb F_2[y]$,
and K\"unneth gives the rank-two formula.

On a cyclic subgroup whose order is divisible by four, the two
restricted degree-one classes are multiples of one exterior
generator, so their product vanishes.  If its order has two-primary
part exactly two, both restrictions vanish: an element of order two
in $(\Z/2^a)^2$ has even coordinates.  A cyclic subgroup of odd
order has no positive mod-two cohomology.
Finally, transfer annihilates positive cohomology for an odd-order
group, because its order is invertible in $\mathbb F_2$.
\end{proof}

If a finite class is the restriction of
$\alpha\in H^2(BE;\mathbb F_2)$, the inclusion $E[2]\hookrightarrow E$
sends $u_i$ to $x_i^2$.  Such a restriction has $b_{11}=0$.
A stabilizer class obtained from an equivariant gerbe on a moduli
space must be computed from that gerbe.  It is not an arbitrary class
of $BE$, and it is not assumed independent of its global restrictions.

\begin{remark}[Square roots and their equivariant choices]
\label{rem:translation-linearization-criterion}
For a specified square-root gerbe on a stabilizer classifying space,
a trivialization exists only when its degree-two obstruction vanishes.
When it does, isomorphism classes of trivializations form a torsor
under $H^1(BH;\mathbb F_2)$.
The degree-two obstruction and the degree-one choice have distinct
roles.  An invariant Picard class does not choose either an
equivariant cocycle or a trivialization.

An equivariant line descends to the quotient \emph{stack} with its
stabilizer representation.  Triviality of that representation is an
additional requirement for descent to a coarse quotient, or for a
specified invariant frame.  Thus a demand that a character vanish
must name its descent target.  These cohomology calculations give
criteria once the actual orientation class is supplied.  They do not
compute that class for the compact Hall moduli.
\end{remark}

\begin{remark}[Characters of odd order]
\label{rem:mod2-not-odd-order}
A sign character detects only two-primary information.  A line
linearization has a character in $\Hom(H,\C^\times)$, which can be
nontrivial for $H$ of odd order.  For example, a generator of
$\Z/3$ can act by a primitive cube root of unity, while
$H^1(B(\Z/3);\mathbb F_2)=0$.  Descent of that line requires its
actual character, not only its sign reduction.
\end{remark}

\subsection{Cusp polarization and chamber transport}
\label{subsec:cusp-polarization-chamber}

\begin{definition}[Cusp order]
\label{def:cusp-polarization}
The standard cusp positivity convention orders $(m,n,-\ell)$
lexicographically.  Inside the positive semigroup
$\Gamma_{\mathrm{eff}}$, retain the Jacobi-active indices
$\Gamma_{\mathrm{act}}$ defined in
\S\ref{subsec:igusa-variables-gram-index}.
A region where every active monomial has absolute value less than
one is $\Imag\tau=\Imag\sigma=t$, $\Imag z=-a$, with $0<a<t$.
Taking $t$ sufficiently large gives the usual cusp region.
\end{definition}

For active indices with $n+m>0$, the inequality
$\ell^2\le4nm+1$ and integrality give $|\ell|\le n+m$.
Thus $t(n+m)-a\ell>0$.  The remaining index $(0,-1,0)$ has
imaginary exponent $a>0$.  Hence each active monomial has absolute
value less than one in this chamber.  This is an order on Fourier
indices, not a positivity theorem for microscopic charges.

\begin{remark}[Arithmetic chamber transport]
\label{rem:chamber-transport}
The index form
\[
\langle\gamma,\gamma'\rangle_{\mathrm{ind}}
=2(nm'+n'm)-\ell\ell'
\]
is $-1/2$ times the root pairing under $\alpha$.
An element of $\Sp_4(\Z)$ transports the automorphic section and its
expansion to a cusp with transported local coordinates.
This arithmetic action does not specify an action on every sheaf
charge or on a stability condition.  Such a geometric action requires
its own functor and charge comparison.  

In the four-dimensional
$N=4$ dyon problem, the electric and magnetic vectors lie in a
rank-$28$ Narain lattice \cite[\S1]{DVV}.  The three quadratic
invariants enter a genus-two generating function.  This does not
identify its Siegel modular group with the full microscopic duality
group or identify the rank-$24$ Mukai lattice with that Narain lattice.
\end{remark}

\begin{remark}[Wrapped objects under a change of chamber]
\label{rem:chamber-dependence-wrapped}
Projection to the fixed curve $E$ determines whether an object is
wrapped.  Changing an automorphic cusp alone does not change that
projection.  A comparison that transports the local and wrapped
categories must specify the functors and prove the corresponding
support statements.
\end{remark}

\begin{definition}[Chamber-compatible Hall comparison]
\label{def:chamber-compatible-hall-descent}
Choose two stability conditions and cusp expansions.
A chamber-compatible comparison consists of maps of their finite
moduli and extension systems, with coherent coefficient and orientation
maps.  Require them to intertwine the charge maps in the chosen
coordinates, preserve the indicated support categories, and commute
with all transition maps.  Wall-crossing identities of scalar
invariants alone do not supply such a comparison.
\end{definition}

The identity
$\overline\Pi_X(\widehat c\star\widehat c')
=\overline\Pi_X(\widehat c)+\overline\Pi_X(\widehat c')$
holds before any chamber is chosen.  It remains a lattice identity
until these categorical maps have been constructed.

\subsection{Constructible coefficients for local operations}
\label{sec:local-protected-observable-algebra}

Deformation and obstruction complexes on stable-sheaf moduli underlie
Donaldson--Thomas theory \cite{DT}.  A Hall construction retains
cohomological coefficients and extension operations before taking
numerical invariants.  Comparing it with an algebra of BPS states,
as in \cite{HM}, requires a map between the specified state spaces
that preserves their products.

For an open $U\subset X$, select finite-type moduli
$\mathfrak M_c(U)$ of the desired compactly supported objects.
Assume they carry oriented critical charts and a constructible
vanishing-cycle complex $\Phi_c$ with compatible descent.  This
requires a specified coefficient formalism for these stacks.
On oriented d-critical schemes the construction is provided by
\cite[Theorem~6.9]{BBDJS2015}.  Applying it to a quotient requires
compatible charts, orientations and equivariance.

\begin{definition}[Protected coefficient complex]
\label{def:local-protected-observable-algebra}
Suppose $R\Gamma_c(\mathfrak M_c(U),\Phi_c)$ is a bounded complex
of finite-dimensional rational vector spaces.  Fix the convention
\[
C^{\mathrm{BM}}(\mathfrak M_c(U);\Phi_c)
=R\Hom_\Q\bigl(R\Gamma_c(\mathfrak M_c(U),\Phi_c),\Q\bigr),
\qquad
\mathcal A_X(U)=\bigoplus_c
C^{\mathrm{BM}}(\mathfrak M_c(U);\Phi_c).
\]
Here the sum runs over the specified geometric charge support.
An algebra structure on $\mathcal A_X(U)$ consists of homogeneous
maps induced by the selected extension spans, with stated Gysin
shifts, pushforward hypotheses and coherent coefficient comparisons.
\end{definition}

For $\Phi_c=\Q$, the convention is the usual Borel--Moore complex
$R\Gamma(\mathfrak M_c(U),\omega_{\mathfrak M_c(U)})$ by Verdier
duality.  The orientation is used to construct $\Phi_c$.  It is a
coherent square-root line on the critical geometry, so it cannot be
tensored directly with a rational constructible complex.  A change
of orientation can instead produce a specified sign local system
in the constructible category.  Orientation also enters the proposed
vanishing-cycle definition of Gopakumar--Vafa invariants
\cite[\S1.2]{MaulikTodaGV}.

For disjoint opens $U_1,\ldots,U_k$, objects on their disjoint union
decompose among the opens.  If $F$ and $G$ have disjoint supports,
the sheaf $R\mathcal Hom(F,G)$ is zero, hence $R\Hom(F,G)=0$.

This assertion does not require compact support.  Compactness and
finite type enter the coefficient functors and pushforwards.
When the chosen moduli and orientations respect this product
decomposition, Thom--Sebastiani and K\"unneth give maps
\[
\mathcal A_X(U_1)\otimes\cdots\otimes\mathcal A_X(U_k)
\longrightarrow\mathcal A_X(U_1\sqcup\cdots\sqcup U_k)
\]
\cite{MasseyST,SaitoMHM}.  Extension to overlapping opens and
collision strata requires further factorization data.
A comparison with the D-module formalism of \cite[\S3.4]{BD}
requires a realization on configuration spaces, compatible with
their diagonals.

The Lie--commutative duality of \cite{FrancisGaitsgoryCKD} concerns
chiral Lie algebras and commutative coalgebras in the chiral tensor
category.  An associative-to-associative comparison
\cite{ChiralBarCobarVolI,ChiralBarCobarVolII} additionally requires
the associative products and compatible collision maps on both sides.
These operations are not supplied by the disjoint-support calculation.

For the unreduced moduli of coherent sheaves and perfect complexes
on a smooth projective Calabi--Yau threefold, orientation data do
exist \cite[Theorem~4.4]{JoyceUpmeier}.  The theorem applies to
$S\times E$ with the chosen trivialization of its canonical bundle.
It supplies isomorphism classes of square roots compatible with
the indicated exact-sequence or triangle operations.  Compatibility
with a chosen scalar and trace reduction, translation quotient,
boundary degeneration and higher flag maps must still be established
for the system considered here.

\begin{remark}[Scope of the coefficient construction]
\label{rem:source-interface-external-assumptions}
The shifted symplectic form, an oriented critical chart, its
constructible complex, and a Hall product are separate structures.
The product needs actual pull-push maps on extension stacks.
Multiplicative orientation requires their compatibility on longer
flags.  Translation reduction requires the common-torsor spans and
equivariant descent described above.  None is constructed by the
scalar superdimension of a coefficient complex.
\end{remark}

\subsection{Orientation and modular comparison}
\label{sec:k3xe-four-obstructions}

Four geometric requirements organize the orientation comparison.
Condition $(D0)$ asks for a D0 degeneration with compatible reduced
complexes, semiregularity cosections, coefficient specialization,
orientation and transition maps.  Condition $(O1)$ asks for a square
root of the reduced determinant line on the elliptic quotient,
including its boundary behavior.  Condition $(O1)^+$ asks for coherent
Weyl lifts and a trivialization of their orientation cocycle.
Condition $(O2_{\mathrm{atlas}})$ asks for compatible type-II wall
charts with a normal skew complex
\([V\xrightarrow{D_u}V^\vee]\) in degrees \([-1,0]\), where
\(V=\mathcal O^2\) and \(\operatorname{pf}(D_u)=u\).
The chart includes a compatible reflection lift on the root line
and a reflection-invariant tangential unit.
The corresponding conditional criteria are stated in
Appendix~\ref{app:obstruction-discharges}.

These orientation requirements do not alone identify a geometric
Pfaffian with an automorphic form.  The finite geometric datum
$(P_{\mathrm{fin}})$ of
Theorem~\ref{thm:pentadic-pfaffian-dirac} also specifies the reduced
skew complexes, Pfaffian sections, Fourier coefficients, line maps,
cusp frames and compatible limits.  Under all those hypotheses, the
line and section comparison gives
\[
\operatorname{Pf}_{\mathrm{prot}}(\mathfrak D_X^{\mathrm{DI}})
=\Delta_5,
\qquad \epsilon_o=\widetilde\nu_{\Delta_5}
\quad\text{on }W^{(2)}(\Lambda_{II}^{2,1}).
\]
Construction of the compact geometric hypotheses remains open.

The scalar normalization is a different statement.  In the primitive
K3 class regime and variable convention of the opening chapter,
\[
Z^X_{\mathrm{OP}}=-\frac1{\chi_{10}^{\mathrm{OP}}}
=-4096\Delta_5^{-2},\qquad
\chi_{10}^{\mathrm{OP}}=(\Delta_5/64)^2.
\]
The reduced stable-pair and Gromov--Witten comparisons underlying
this identity have the hypotheses stated there
\cite{OberdieckReducedPT,OPand,OB}.
The reduced Donaldson--Thomas theory uses weighted Euler
characteristics after quotienting the Hilbert scheme by elliptic
translation \cite{BRYAN}.

The normalized scalar
$Z_{\mathrm{BPS}}^{K3\times E}:=-Z^X_{\mathrm{OP}}/4096$
is $\Delta_5^{-2}$.  Calling it a trace of an independently
constructed geometric state space requires the state comparison.
The inverse square loses the sign character of $\Delta_5$.
The formal set of twelve signs has $2^{12}=64^2$ elements.
A geometric orbit of that size requires its own bijection.
The numerical equality supplies no map from an orientation line
to an automorphic line.

\begin{remark}[Numerical invariants and comparison coefficients]
\label{rem:kappa-tuple-compatibility}
The tuple
\[
(\chi(\mathcal O_X),\sum_q(-1)^qh^{0,q}(X),
\rk\Gamma_{\mathrm{gram}},\operatorname{wt}\Delta_5,
\chi_{\mathrm{top}}(S))=(0,0,3,5,24)
\]
records five defined numerical invariants.  Replacing its third entry
by a genus-one chiral or Heisenberg coefficient requires a chosen
complex, trace normalization and a comparison with the Gram rank
\cite{CYQGVolIII}.
A matching number does not supply this comparison.  Nor does the
scalar Igusa identity construct a gravitational path integral on
the compact target.  A proposed local holomorphic-topological model
needs a global gluing map, with its obstruction class and coefficient
comparison, before it can yield that interpretation
\cite{MixedHTSCompanion}.
\end{remark}

\subsection{Fixed notation and the remaining maps}
\label{sec:k3xe-notation-summary}

The cohomological charge group is $\Gamma_X^{\mathrm{form}}$, its
algebraic subgroup is $\Gamma_X^{\mathrm{alg}}$, and geometric support
is selected inside it.  The map $\Pi_X$ is quadratic.
The split extension $\widehat\Gamma_X$ carries the additive map
$\overline\Pi_X$.  The hybrid parameter space records local and
wrapped families, while their operations require the extension spans
and coefficients just specified.

At each height, the proposed Dirac--Igusa system consists of these
geometric data together with its Hall and primitive maps, Koszul
comparison, Pfaffian line and section.  A pro-object exists from a
coherent cofiltered system of such stages.  The construction problem
is to supply the compact stages, the comparison maps and their
transition identities.  The lattice extension, support lemma and
finite regrading results solve their stated parts without replacing
that problem.
